\chapter*{Introduction}
\addcontentsline{toc}{chapter}{Introduction}
\markboth{Introduction}{Introduction}

\section*{Working definition of periodic mathematics}

\begin{definition}[Periodic mathematics (working definition)]
Periodic mathematics is a \textbf{research programme} proposed in this book. It takes the \textbf{symmetric structures generated by group actions} as a unified language, \textbf{translation, scaling and rotation} as the candidate classification of the basic generating transformations, and the \textbf{closed orbit} as the core geometric picture, and it attempts a unified mathematical description of the periodic and fractal structures found in different disciplines. Its objects of study are systems possessing a state space, evolution rules and a symmetric structure. Periodic mathematics does not assert that all systems are periodic; rather, it is devoted to providing a unified diagnostic and descriptive language for \textbf{systems in which a periodic structure can be defined}. \textbf{It must be stressed} that the above formulation is a working definition and a research programme, not a formal mathematical definition; its rigorous version must be established chapter by chapter, with conditions specified one by one.
\end{definition}

\section*{Philosophical inspiration and cognitive method}

The saying from the \emph{Tao Te Ching}, `The Dao gives birth to the one, the one to the two, the two to the three, and the three to the myriad things', and the saying from the \emph{Yin Fu Ching}, `Observe the way of heaven; hold fast to the movement of heaven', provide the philosophical starting point and the \textbf{inspiration} for the cognitive method of this book. The former describes a generative sequence: starting from a single origin, and passing through the ladder of differentiation, it produces diversity. In periodic mathematics, this sequence \textbf{can be likened to} the ladder of number-system extensions $\mathbb{N} \to \mathbb{Z} \to \mathbb{Q} \to \mathbb{R} \to \mathbb{C} \to \mathbb{H}$---each step of extension introduces a new inverse operation, and each inverse operation introduces a new symmetry. This analogy is heuristic, not a mathematical theorem. The latter suggests a methodological path: observe invariance, abstract conservation laws, discover symmetries, identify closed orbits. \textbf{Observing the way of heaven} may be understood as identifying the conserved quantities of a system, and \textbf{holding fast to the movement of heaven} as using conservation laws to predict the future states of the system.

\section*{Leibniz's universal characteristic and the mathematical foundations of the \emph{I Ching}}

After reading the hexagram images of the \emph{I Ching}, Leibniz wrote: `This discovery convinces me that Fuxi was not only the founder of Chinese philosophy, but also one of the most ancient sages of all humanity. The system of symbols he left behind contains a universal writing---a character capable of expressing all thought.' \textbf{His concrete discovery}---the strict periodicity of the digits in each place of the binary expansion of the natural numbers, with periods $2$, $4$ and $2^{k+1}$ for the $0$th, $1$st and $k$th places---together with the examination of the historical sources as to whether he named it \emph{Periodus numerorum}, has already been given in the \textbf{Preface}, and is not repeated here.

A further extension in the direction of the \emph{I Ching} is this: the generative structure of the hexagram system corresponds to the approximation of the circle by binary fractions (see the \textbf{Preface}), and the saying in the \emph{I Ching}, `The \emph{I} is modelled on heaven and earth', can \textbf{partly} acquire mathematical meaning in periodic mathematics: there exists an isomorphism between the hexagram system and certain periodic structures. `Those of the same sound respond to each other, those of the same vapour seek each other' can be likened to the resonance of frequencies and the attraction of phases; `There is no plain without a slope, no going without a return' can be likened to the necessity of closed orbits; `Things cluster by their kind, creatures divide by their groups' can be likened to the classification of orbits under a group action. These analogies are \textbf{heuristic}, and their mathematical content must be defined one by one in the concrete chapters.

The periodic convolution of the Heavenly Stems ($T=10$) and the Earthly Branches ($T=12$) generates the sexagenary cycle, and the nesting of the Eight Trigrams ($T=8$) with the Eight Trigrams ($T=8$) generates the Sixty-Four Hexagrams---the epicycle nesting rule unifies the process by which small periods generate large periods: \textbf{with each step of the outer circle, the inner circle rotates a full turn}. Yin-Yang multiplied by the Five Phases gives the ten Heavenly Stems, the Four Symbols multiplied by the Three Powers gives the twelve Earthly Branches, and the Eight Trigrams multiplied by the Eight Trigrams gives the Sixty-Four Hexagrams. \textbf{In this book these correspondences are presented as algebraic isomorphisms or as equivalences of group actions, and not as assertions that the ancient texts themselves were already expressed in the language of modern group representation theory}.

\section*{Why translation, scaling and rotation}

Periodic mathematics establishes translation, scaling and rotation as the three basic generating transformations, not on the basis of empirical induction, but on the basis of \textbf{motivation} at the following three levels. \textbf{It must be explained} that the three points below are \textbf{reasons for choosing these three}, and not the mathematical theorem that `any continuous transformation can be decomposed as a composition of the three'; the latter is a stronger proposition that requires case-by-case verification.

\begin{enumerate}[label=(\roman*)]
    \item \textbf{Algebraic level}: the three transformations correspond respectively to the three most basic continuous groups---translation to the additive group $(\mathbb{R}^n,+)$, scaling to the multiplicative group $(\mathbb{R}^+,\times)$, and rotation to the circle group $U(1)$. \textbf{At a suitable level of abstraction}, many continuous transformations that preserve some intrinsic structure can be decomposed as a composition of these three groups; but this decomposition does not hold for all continuous groups, and must be decided case by case.
    \item \textbf{Geometric level}: translation changes position without changing orientation, scaling changes scale without changing shape, and rotation changes orientation without changing position. \textbf{Within the framework of Euclidean and similarity geometry}, these three operations exactly exhaust the continuous transformations that leave the intrinsic structure of a system unchanged.
    \item \textbf{Number-system level}: each step of number-system extension introduces a new inverse operation---subtraction introduces the translational direction (centre of symmetry $0$), division introduces the proportional direction (centre of symmetry $1$), and complex roots introduce the rotational direction (centre of symmetry $i$). This correspondence is \textbf{heuristic}, and its rigorous mathematical content belongs to the study of field extensions and automorphism groups.
\end{enumerate}

\section*{Gradient, curl, divergence and the correspondence with the three symmetries}

In vector analysis, the gradient $\nabla$, the curl $\nabla\times$ and the divergence $\nabla\cdot$ are three basic differential operators. Their individual structures are as follows (the rigorous definitions, hypotheses and boundary conditions of the three are given in the `Vector Analysis' chapter; this book does not place them side by side with the three basic transformations as a correspondence):

\begin{enumerate}[leftmargin=4em, label={(\arabic*)}]
    \item \textbf{Gradient}: the gradient maps a scalar field to a vector field and points in the direction in which the function increases fastest, its magnitude being equal to the maximal directional derivative (this requires $f \in C^1$ and a non-vanishing gradient at that point).
    \item \textbf{Divergence}: the divergence maps a vector field to a scalar field and measures the net outward flux per unit volume; in a Hamiltonian system the vanishing of the divergence is equivalent to the conservation of phase-space volume (Liouville's theorem).
    \item \textbf{Curl}: the curl maps a vector field to a vector field and characterises the local rotational tendency at each point, its magnitude being equal to twice the local angular velocity (this requires $\mathbf{F} \in C^1$).
\end{enumerate}

\textbf{It must be stressed} that the most standard structure of the divergence is that of local sources and sinks, the rate of volume change, the continuity equation and Liouville-type results, whereas the generator of scaling at the level of continuous groups is $x\,d/dx$, an object of another level; the relation between the two may be studied, but they cannot be placed side by side as a correspondence law before a rigorous bridge has been built. This book therefore does not state `which operator corresponds to which transformation' as a proposition.

\section*{Why conservation laws are needed}

Conservation laws are not an appended physical condition in periodic mathematics, but the \textbf{logical precondition for the stable existence of periodic structure in conservative systems}. \textbf{It must be stressed} that the following arguments apply to conservative systems; for dissipative systems, one must introduce the notion of a generalised conserved quantity or of an attractor, which will be discussed in later chapters.

\begin{enumerate}[label=(\roman*)]
    \item \textbf{Reversibility of transformations}: a transformation can be applied repeatedly without destroying the structure of the system, \textbf{in a conservative system}, if and only if there exists at least one quantity that is strictly invariant under that transformation. A transformation without a conserved quantity is typically irreversible within the conservative framework---the closure of the orbit is then difficult to guarantee logically.
    \item \textbf{Constraint of orbits}: the mathematical essence of a conservation law is to constrain the evolution of the system to an invariant submanifold. Without such a constraint, the evolution of the system is free drift, which \textbf{generally} produces no recognisable regular orbit. \textbf{In dissipative systems, however, an attractor can take the place of a conserved quantity and play the constraining role}.
    \item \textbf{Observable regularity}: a conserved quantity is an identity across time---it establishes a decidable connection between the `state of the past' and the `state of the present'.
\end{enumerate}

Viewed from the algebraic framework established in the `Periodic Foundations' chapter, a periodic system can perfectly well possess only trivial invariants (for example a finite cyclic permutation: the state set is finite, the transformation is of finite order, every state is periodic, and yet a non-trivial conserved quantity need not exist). Therefore \textbf{a conserved quantity is neither a necessary condition nor, generally, a sufficient condition for periodicity}. The fundamental principle of this book is stated as follows:

\textbf{In a conservative system, a conserved quantity can provide an invariant constraint on the orbit, and thereby furnish an important condition for the existence and stability of periodic structure; but a conserved quantity is not a sufficient condition for periodicity, nor, generally, a necessary one. Within the framework of this book, the algebraic hinge that directly produces periods is the finite-order action:}\[
\text{finite-order action} \;\Longrightarrow\; \text{uniformly closed orbit} \;\Longrightarrow\; \text{period}.
\]

The proper role of conservation and symmetry within this is that of \textbf{auxiliary tools for discovering, constraining and explaining periodic structure}, and not of indispensable links in the logical chain.

\section*{Why group representation theory and the manifold framework are needed}

\begin{enumerate}[label=(\roman*)]
    \item \textbf{Group representation theory answers `what the transformation is'}: group representation theory provides the algebraic structure of transformations---does the set of transformations form a group? Do the group elements have finite order? What is the structure of the orbit?
    \item \textbf{The tangent space answers `what the instantaneous direction of the transformation is'}: the tangent space captures the infinitesimal direction of change of the system in its current state. The tangent vector at each point of a closed orbit constitutes the instantaneous velocity field of that periodic motion. Whether the vector field in the tangent space is integrable determines whether the orbit can close within finite time.
    \item \textbf{The normal space answers `what constrains the transformation'}: the normal space describes the constrained directions experienced by the system during its evolution. Conservation laws appear geometrically as constraint hypersurfaces in the normal space---the system is confined to move on this hypersurface and cannot escape. The dimension of the normal space is equal to the number of conserved quantities, and the codimension of the normal space is equal to the degrees of freedom of the system.
    \item \textbf{The angle space answers `how far the transformation has progressed'}: the angle space records the phase position of the system on the closed orbit. For a system of period $T$, the angle space is a parametrisation by the circle group $U(1)$ or by the finite cyclic group $\mathbb{Z}_T$. The introduction of the angle space turns `phase' into an independent geometric dimension, and not merely a function of time. \textbf{It should be noted} that the specific type of the angle space depends on the fundamental group and the homotopy class of the orbit; not every closed orbit automatically possesses the $U(1)$ or $\mathbb{Z}_T$ structure, and its rigorous classification requires topological conditions.
    \item \textbf{The manifold framework answers `what the global structure is'}: a manifold is the global integration of the tangent space, the normal space and the angle space. The complete state space of a periodic system is a manifold; a closed orbit is a closed curve on this manifold, the tangent space is the tangent plane of the manifold at each point of the orbit, the normal space is the orthogonal complement of the tangent plane, and the angle space is the parametrisation of the orbit itself.
\end{enumerate}

\textbf{Only the combination of all five yields a complete determination}: group representation theory provides the algebraic possibility, and the manifold framework provides the geometric stage. Only by uniting the group action, the tangent space, the normal space and the angle space within the manifold framework can one answer simultaneously: algebraically, do the group elements have finite order? geometrically, does the orbit close? in the tangent direction, does an integrable vector field exist? in the normal direction, does a conserved constraint exist? in the angle direction, is the phase single-valued? \textbf{This framework is a research programme}, and its specific applicability conditions need to be established case by case in the various chapters.

\section*{Formalisation of the manifold framework}

Let $M$ be the state-space manifold of the system, $\gamma: [0,T] \to M$ be a closed orbit, and $p = \gamma(t)$ be a point on the orbit. Then:

\begin{enumerate}[leftmargin=4em, label={(\arabic*)}]
    \item \textbf{Tangent space} $T_p M$: the one-dimensional subspace spanned by the tangent vector $\dot{\gamma}(t)$ of the orbit at the point $p$, describing the instantaneous direction of motion of the system;
    \item \textbf{Normal space} $N_p M$: the orthogonal complement of the tangent space in $T_p M$, describing the constrained directions experienced by the system. The gradient $\nabla Q(p)$ of a conserved quantity $Q$ is a vector in the normal space;
    \item \textbf{Angle space} $A_p$: the parametrisation space of the orbit $\gamma$; for a periodic orbit, $A_p \cong U(1)$ or $A_p \cong \mathbb{Z}_T$, describing the phase position of the system on the orbit. \textbf{Conditions are required}: when the element of the fundamental group of the orbit is of infinite order and the orbit is contractible, the angle space is naturally $U(1)$; when the element of the fundamental group is of finite order, the angle space is $\mathbb{Z}_T$. The general case requires separate discussion.
    \item \textbf{Tangent bundle} $TM = \bigcup_{p \in M} T_p M$: the collection of all tangent spaces, describing the directions of motion of the system in all possible states;
    \item \textbf{Normal bundle} $NM = \bigcup_{p \in M} N_p M$: the collection of all normal spaces, describing the constrained directions of the system in all possible states;
    \item \textbf{Angle bundle} $AM = \bigcup_{p \in M} A_p$: the collection of all angle spaces, describing the phase positions of the system in all possible states.
\end{enumerate}

\section*{Correspondence between conservation laws and the normal space}

Conservation laws acquire a geometric interpretation within the manifold framework: \textbf{under suitable non-degeneracy conditions}, a conserved quantity $Q$ corresponds to a constraint hypersurface $Q^{-1}(c)$ in the normal space. The system is confined to move on this hypersurface, and the orbit $\gamma$ is a closed curve on the constraint hypersurface. The vectors in the tangent space $T_p M$ are orthogonal to $\nabla Q(p)$, which means that the conserved quantity does not change along the direction of motion. The dimension of the normal space is equal to the number of independent conserved quantities, and the codimension of the normal space is equal to the effective degrees of freedom of the system. \textbf{When the number of independent conserved quantities is equal to the degrees of freedom of the system}, the orbit is completely constrained to be a closed curve---this is a geometric hallmark of a completely integrable system; \textbf{but the converse does not hold}: the existence of a closed orbit does not require the number of independent conserved quantities to be equal to the degrees of freedom.

\section*{Correspondence between the angle space and the phase operator}

The angle space is the geometric realisation of the phase operator. The $k$th phase of period $T$ corresponds to the complex number $\mathcal{O}_k^{(T)} = e^{i2\pi k/T}$, which appears in the angle space as a uniformly sampled point on the orbit; its three actions (identity, antipode, rotation by $2\pi k/T$) correspond one by one to the algebraic description of the phase operator in the Preface, and are not repeated here. The topological structure of the angle space determines the algebraic structure of the phase operator: when the angle space is the continuous circle $U(1)$, the phase operator is a continuous rotation group; when the angle space is the finite cyclic group $\mathbb{Z}_T$, the phase operator is a finite rotation group. \textbf{It should be noted} that the above correspondence presupposes that the angle space is well defined, and the rigorous construction of the angle space depends on the homotopy class and the fundamental group of the orbit; not every closed orbit automatically satisfies this.

\section*{A preliminary definition of periodic convolution}

Traditional harmonic analysis mainly decomposes a signal into a sum of sinusoidal waves of different frequencies, and does not necessarily preserve the interaction information between amplitude and phase. A core research programme proposed in this book is periodic convolution, whose motivation is to bring amplitude (scaling) and phase (rotation) into the series operations as well.

\begin{definition}[Periodic convolution: interaction tensor and composition map (preliminary framework)]
Let two periodic systems be represented by the epicycle series
\[
f(x) = \sum_{n} A_n e^{i(\omega_n x + \phi_n)}, \quad g(x) = \sum_{m} B_m e^{i(\nu_m x + \psi_m)}
\]
(under suitable convergence conditions).

\textbf{First layer (interaction tensor)}: define the pairwise interaction tensor between the two families of epicycle components
\[
C_{n,m} = A_n B_m \operatorname{Similarity}^{(d)}(\theta_n, \theta_m),
\]
where $\theta_n = (\omega_n, \phi_n)$ is the phase point of $f$, $\theta_m = (\nu_m, \psi_m)$ is the phase point of $g$, and $\operatorname{Similarity}^{(d)}$ is the $d$-dimensional phase-similarity kernel.

\textbf{Second layer (composition map)}: define the composition map from the interaction tensor to the function space
\[
S(C)(x) = \sum_{n,m} C_{n,m} e^{i(\Omega_{n,m} x + \Phi_{n,m})},
\]
where $\Omega_{n,m}$ is the frequency-composition operator and $\Phi_{n,m}$ is the phase-composition operator.

\textbf{Periodic convolution} (preliminary) is defined as the composition of the interaction tensor and the composition map:
\[
(f \circledast g)(x) \;=\; \mathcal{K}^{(d)}\{f,g\} \;=\; S(C)(x).
\]
\end{definition}

\textbf{Notation}: throughout the main text of this book the periodic-convolution operator is written $\circledast$; when the order $d$ of the phase-similarity kernel $\operatorname{Similarity}^{(d)}$ has to be displayed, or when a window kernel $W$ is introduced, the same operation is written $\mathcal{K}^{(d)}\{f,g\}$---this is the form used on the cover---and in the applied chapters the shorthand $A_{n,m} := A_n B_m$ is occasionally employed.

\textbf{It must be emphasised} that the above definition is a \textbf{research programme}, and that for it to hold rigorously the following conditions must be specified: (i) the series participating in the convolution converge under a suitable topology; (ii) the phase-similarity kernel $\operatorname{Similarity}^{(d)}$ satisfies symmetry, boundedness and normalisation conditions; (iii) the frequency-composition operator and the phase-composition operator are explicitly defined and satisfy associativity or distributivity. These conditions will be given one by one in later chapters or listed as open problems. \textbf{Until the rigorous definition has been completed, the above framework should be regarded as a heuristic programme rather than an established mathematical object.}

Periodic convolution takes amplitude (scaling), frequency (translation), phase (rotation) and period (fixed length) as its four operator parameters, and preserves the interactions of all frequency pairs and all period pairs. This means that the interaction between two periodic systems, \textbf{within the framework of periodic convolution}, depends not only on their respective frequencies, but also on the ratio of their amplitudes, the difference of their phases, and the common-multiple relation between their periods. This book will attempt to argue that the classical integral transforms---the Fourier transform, the wavelet transform, the Mellin transform and the Laplace transform---\textbf{can be given a unified description within the framework of group representation and harmonic analysis under suitable group actions, representations and function spaces}; \textbf{this is more cautious than saying that `they are essentially one and the same thing'}---since wavelet analysis involves the affine group, translation plus scaling, representations, continuous and discrete wavelets, admissibility and specific measures, so that the precise unified form requires case-by-case verification. The rigorous definition and properties of periodic convolution will be established in later chapters.

\section*{An information ontology and the observer framework}

\textbf{The information ontology proposed in this section is a philosophical position of this book, not a mathematical theorem.} Its mathematical content lies solely in the fact that the scale and the sampling rate of the observer determine the range of information that can be mathematised by the observer.

Let the composite vector of the temporal system be $\mathbf{V}_T$, the composite vector of the spatial system be $\mathbf{V}_S$, and the composite vector of the observer system be $\mathbf{V}_P$. Then
\begin{align}
\|\mathbf{V}_T\|^2 &= E \quad (\text{energy}), \label{eq:energy}\\
\|\mathbf{V}_S\|^2 &= m \quad (\text{mass}), \label{eq:mass}\\
\|\mathbf{V}_P\|^2 &= \sigma \quad (\text{scale}). \label{eq:scale}
\end{align}
Energy, mass and scale---here \textbf{interpreted} as the squared moduli of the composite vectors of three independent information-flux systems. The \textbf{mathematical content} of this formalisation will be given in later chapters; its \textbf{ontological interpretation} is philosophical and does not belong to the category of mathematical theorems.

The observer itself is also part of the information structure. The scale $\sigma$ of the observer and the sampling rate $R(r,\sigma)$ directly determine the range of information that can be mathematised. When the number of sampled points $N(r,\sigma) < 1$, the information is completely invisible to the observer; when the sampling rate $R$ is insufficient, high-frequency information is aliased into low-frequency artefacts. \textbf{The boundary of the subset of information that has been mathematised is a function of the observer's own sampling limit}---this is a structural statement about the relation between the observer and information, whose rigorous form is given in the `Discrete Sampling Approximation, Congruence Angles, and Intercalation' chapter, \cref{def:dynamic_sampling}, \cref{thm:invisibility_criterion}.

The choice of the observer's sampling window is a key technical step in the `mathematisation of information'. When the continuous period $T_f$ being observed and the sampling interval $\Delta x$ of the observer are incommensurable, the choice of the sampling window length directly determines the quality of the spectrum that the observer can resolve. The \textbf{intercalation method} introduced in this book (see the Preface) provides the observer with a window-selection mechanism that has a rigorous optimality criterion: by selecting the continued-fraction convergents $p_j/q_j$ of $T_f/\Delta x$, the observer can obtain the minimal phase error and the minimal spectral leakage on the window length $T = p_j$. This mechanism transforms the observer framework from a philosophical statement into an operable algorithmic criterion. The canonical definitions of the sampling interval and the sampling rate are given in the `Discrete Sampling Approximation, Congruence Angles, and Intercalation' chapter, \cref{def:uniform_sampling_parallel,def:nyquist_sampling_parallel}.

\section*{The duality of information}

\begin{enumerate}[leftmargin=4em, label={(\arabic*)}]
    \item \textbf{Independent existence}: information can exist independently as an identified unit---a group element, a function value, an epicycle radius, a phase vector.
    \item \textbf{Interrelation}: information can likewise manifest itself as the relation between entities---group operations, function composition, the nested superposition of epicycles, the periodic-convolution projection of phase vectors.
\end{enumerate}

The \textbf{reification} of relations can form the foundation of a hierarchical information structure: first-order information (entities), second-order information (relations between entities), third-order information (relations between relations). \textbf{Quotient groups, functors and natural transformations can be viewed as several mathematical realisations of this hierarchical leap in information}---this correspondence is heuristic, and its rigorous form depends on the specific category-theoretic construction.

\section*{Evolution model: the arising, abiding and ceasing of conservation laws}

The time derivative of a conserved quantity
\[
\frac{dQ}{dt} = \{Q, \mathcal{H}\} + \frac{\partial Q}{\partial t}
\]
explicitly distinguishes internal dynamics from external driving. The three stages of generation ($\nabla Q \to 0$), evolution ($\nabla Q = 0$) and breaking ($\nabla Q \neq 0$) \textbf{can be given a rigorous formulation within a suitable symplectic-geometric framework}.

\section*{Research aims and methods}

The research aim of periodic mathematics is to establish a unified mathematical framework capable of describing the \textbf{common symmetric structure} of different dynamical systems. We are not concerned with the specific physical, biological or social mechanisms of each system, but with their \textbf{common mathematical structure}.

Periodic mathematics accepts all the standard constructions of modern mathematics and does not redefine existing concepts. Its concrete work is:
\begin{enumerate}[leftmargin=4em, label={(\arabic*)}]
    \item To reveal symmetry and directionality as the foundations of periodicity---the three fundamental symmetries ($0$, $1$, $i$) correspond respectively to translation, scaling and rotation;
    \item To give a rigorous definition of periodicity in the language of group representation theory---a periodic function is an invariant function under a group action, and an epicycle is the geometric realisation of rotational symmetry;
    \item To reveal the structural correspondence between the extension of number systems and the growth of symmetry groups;
    \item To incorporate the observer framework of dynamic sampling in a unified manner;
    \item To incorporate integral transforms in a unified manner into the harmonic-analysis framework on groups with periodic orbits \textbf{under suitable conditions};
    \item To use periodic convolution for a unified description of cross-period phase operations---phase as a relational operator, epicycles nested to generate large periods, the sexagenary cycle being the periodic-convolution composition of the Heavenly Stems and the Earthly Branches;
    \item To establish the \textbf{intercalation method} as a general approximation tool for incommensurable periodic systems---with the best continued-fraction approximation as its theoretical foundation, the bounded-intercalation-error theorem as its rigorous support, and the intercalation functor as its category-theoretic realisation, so as to reduce the spectral leakage in periodic analysis from uncontrollable to bounded.
\end{enumerate}

\section*{Relationship with existing mathematics}

Periodic mathematics does not seek to replace or revise existing disciplines such as group representation theory, harmonic analysis and representation theory, but rather to reorganise, from the perspective of `period', those parts of these disciplines that are related to periodic structure. The central insight underlying this unity is:
\begin{center}
\textbf{The gradient, the curl and the divergence are rigorously defined differential operators in vector analysis, each with its own differentiability hypotheses and domain conditions; this book asserts no correspondence between them and the three basic transformations, and any relation between the two triples must be established case by case.}\\
\textbf{If suitable group actions, representations and function spaces are chosen, different integral transforms can be given a unified description within the framework of group representation and harmonic analysis.}\\
\textbf{Energy, mass and scale can be interpreted as the observer's projection of the squared moduli of the composite vectors of three flux systems---this is a research programme, not a proven theorem.}\\
\textbf{Intercalation is a bounded-error method for approximating irrational period ratios by integer sequences, and its optimality is given by continued-fraction convergents---this is an independent contribution of this book to the methodology of periodic analysis.}
\end{center}

This manuscript is positioned as a \textbf{research programme}: it provides a new way of organising material and a new angle from which to pose questions, and whether it ultimately holds depends on whether non-trivial new mathematical results can be produced in the future.

\section*{Research principles and steps}

\begin{enumerate}[leftmargin=5em, label={\textbf{Principle \arabic*}}]
    \item \textbf{Structure takes priority over mechanism}. Focus on the common mathematical structure among systems, rather than on system-specific causal mechanisms.
    \item \textbf{Finite-order action as the core}. Within the framework of this book, the algebraic hinge that directly produces periods is the finite-order action: finite order $\Rightarrow$ uniformly closed orbit $\Rightarrow$ period. Conservation and symmetry are not universal necessary conditions, but auxiliary tools for discovering and constraining periodic structure.
    \item \textbf{Diagnose rather than assert}. Provide tools to measure the degree of closure of a system under the various symmetric transformations.
    \item \textbf{Periods arise and cease}. The maintenance of conservation laws safeguards the evolution of periods; the intervention of dissipation or external driving leads to the cessation of periods.
\end{enumerate}

\begin{enumerate}[leftmargin=5em, label={\textbf{Step \arabic*}}]
    \item \textbf{Construction of the mathematical core}. Rigorously define the mathematical propositions and theorems leading from group actions and orbit closure to conserved quantities.
    \item \textbf{Infinitesimal probes}. Use the gradient, the curl and the divergence as candidate tools for the differential expression of the three symmetries.
    \item \textbf{Categorical unification}. Use category theory to formalise different types of periodic structure.
    \item \textbf{Interdisciplinary diagnosis}. Use this set of tools to pinpoint precisely why a system possesses or lacks a certain periodic structure.
\end{enumerate}

\section*{Principal results of this book}

\textbf{This section lists the results that this book actually establishes.} Each is stated at the strength at which it is proved, with the chapter in which it is located given in brackets; the heuristic pictures and the open problems of the Preface and of the preceding sections are \textbf{deliberately not repeated} here.

\begin{enumerate}[leftmargin=4em, label={(\arabic*)}]
    \item \textbf{The two lines} (Chapter 1, `Two lines: iteration generates operations and transformations'). At the level of parameters, iteration organises lower-order operations into higher-order operations, whilst the parameters of a repeatedly applied transformation ascend the same hierarchy. The \textbf{operation line} is addition $\to$ multiplication $\to$ power (the hyperoperation hierarchy); the \textbf{transformation line} is translation $\to$ scaling $\to$ rotation; the relation between the two lines is a correspondence at the level of parameters, given in the form of a definition and a proposition, and not as an assertion that the two are one and the same thing. The three generating transformations of this book thereby acquire a generative derivation, and not merely an enumeration.
    \item \textbf{The algebraic hinge of periodicity} (Chapters 1 and 13--15). Within the framework of this book, the algebraic hinge that directly produces periods is the \textbf{finite-order action}: finite order $\Rightarrow$ uniformly closed orbit $\Rightarrow$ period. The period set of a periodic function and the structure of the group with periodic orbits are developed in the `Periodic Functions and Groups with Periodic Orbits' chapter; conservation and symmetry enter as auxiliary tools for discovering and constraining periodic structure, and not as indispensable links in the deductive chain.
    \item \textbf{Bounded intercalation error and continued-fraction optimality} (Chapter 21). For an incommensurable period ratio $\rho$, the lower, the upper and the balanced intercalation sequences all satisfy $N_m \in \{\lfloor\rho\rfloor, \lceil\rho\rceil\}$, whilst their cumulative errors remain bounded---lying in $(-1,0]$, $[0,1)$ and $[-1/2,1/2]$ respectively; and the optimal choice of the intercalation window is given by the continued-fraction convergents $p_j/q_j$ of $\rho$. The intercalation method is presented in this book as an independent contribution to the methodology of periodic analysis: it reduces the cost of incommensurability from uncontrollable to bounded.
    \item \textbf{Congruence angles and the generalised Chinese remainder theorem} (Chapter 21). A \textbf{congruence angle} converts the compatibility condition of several periods into an angular condition on the circle group, and generalises the Chinese remainder theorem from the coprime case to a form in which coupled multi-period constraints can be stated, combined and solved; intercalation is thereby upgraded from the suppression of leakage for a single period to the joint suppression of several periods.
    \item \textbf{The two-layer structure of periodic convolution} (Chapters 6, 23 and 45). Periodic convolution is defined as the composition of an interaction tensor and a composition map, $(f\circledast g)(x) = \mathcal{K}^{(d)}\{f,g\} = S(C)(x)$, so that amplitude (scaling) and phase (rotation) enter the series operation; its functoriality on the category of translation periods is stated in the form of a theorem.
    \item \textbf{The epicycle activation function and function approximation} (Chapters 42 and 45). On the basis of the epicycle nesting generation rule---with each step of the outer circle, the inner circle rotates a full turn---the epicycle activation function is defined, and a corresponding function-approximation result is established. It is here that the research programme of this book acquires a computable concrete form.
\end{enumerate}

\textbf{What this book does not claim}: the above are the results that this book establishes. The stronger propositions of the research programme---that any continuous transformation preserving some intrinsic structure can be decomposed into the three generating transformations, that `which operator corresponds to which transformation' constitutes a correspondence law, that a general periodic function is a superposition of finitely many epicycles---\textbf{are not among the claims of this book}, and are carried in the main text only in the form of a research programme or an open problem. What must be distinguished alongside this is the following: within the framework of \textbf{one-parameter families of conformal transformations of the complex plane}, a framework with explicit hypotheses, the completeness of the three basic transformations is a conclusion proved rigorously in section~\cref{sec:completeness_three_transformations} of Chapter 1; what is not claimed above is its generalisation to `an arbitrary continuous transformation group of arbitrary intrinsic structure'. The problems that remain unsolved are stated together in Chapter 46.

\section*{Research structure and interdisciplinary applications}

This book is divided into three parts, together with a closing chapter:
\begin{enumerate}[leftmargin=6em, label={\textbf{Part \arabic*}}]
    \item \textbf{Mathematical foundations and generators} (Chapters 1 to 10). From periodic foundations, number systems, scalars and vectors, functions and limits, real analysis, series and convolution, complex analysis, quaternion analysis, differential equations and analysis, to linear algebra.
    \item \textbf{Core structures, tools and methods} (Chapters 11 to 24). From the foundations of manifolds and fibre bundles, periodic manifolds, periodic categories, probability and distributions, and vector analysis, through epicycle groups and episphere groups, propagation and evolution, groups and group functions, higher-order groups and higher-order functions, and periodic functions and groups with periodic orbits, to discrete sampling approximation (with congruence angles) and intercalation, harmonic analysis, time series analysis, and multi-period coupled systems.
    \item \textbf{Interdisciplinary applications and diagnosis} (Chapters 25 to 45). The periodic structure of probabilistic, cryptographic, astronomical, calendrical, geographic, \emph{I Ching}, electromagnetic, audio, mechanical, circuit-and-chip, biological, psychological, social, climatic, economic and political systems, together with data analysis, agents, the epicycle activation function, algorithmic systems, and modelling from data to epicycles.
\end{enumerate}
The closing \textbf{Chapter 46}, `Summary and Outlook', gathers the three main threads of the book and states the problems that remain open.

The theoretical necessity of interdisciplinary applications lies in this: the value of a unified framework depends on whether it can withstand testing in multiple independent domains. Periodic mathematics not only describes known periodic systems, but also diagnoses those systems that `seem to have a period yet need not satisfy the closure condition'. Planetary orbits, crystal translations, genetic oscillations, the Koch snowflake, Leibniz's binary-sequence periods and the hexagram system of the \emph{I Ching}---\textbf{at a suitable level of abstraction}---share the same structure of group representation theory and the same logic of conservation laws.

In each application chapter, periodic mathematics first asks: \textbf{does this system possess an identifiable symmetry group? Does a corresponding conserved quantity exist? Is the evolution orbit in its state space closed?} (The philosophical statement of the three criteria, and the treatment of the three cases `satisfied', `partly satisfied' and `not satisfied', are given in the \textbf{Preface} under `The stance and limits of this book'.)

\section*{Chapter structure and reading guide}

The conventions and labelling of the chapters of this book (the order of statement `definition $\to$ proposition $\to$ \dots $\to$ remark', and the rules governing the use of the three labels `conjecture / analogy / research programme') are given in the \textbf{Preface} under `Conventions and notation'. The relations among the chapters can be summarised as three main threads:

\begin{enumerate}[leftmargin=4em, label={(\arabic*)}]
    \item \textbf{Operator thread} (gradient $\nabla$ $\to$ curl $\nabla\times$ $\to$ divergence $\nabla\cdot$). Starting from the `Periodic Foundations' chapter, passing through the `Vector Analysis' chapter, and running through the `Periodic Functions and Groups with Periodic Orbits' chapter and the `Harmonic Analysis' chapter, it establishes one by one the differentiability hypotheses and the domain of applicability of each of the three operators. \textbf{This thread asserts no correspondence between the three operators and the three basic transformations}; any relation between the two triples is discussed case by case, and the `Vector Analysis' chapter states explicitly that no correspondence is asserted between them.
    \item \textbf{Number-system thread} ($\mathbb{N} \to \mathbb{Z} \to \mathbb{Q} \to \mathbb{R} \to \mathbb{C} \to \mathbb{H}$). Reveals how periodic structure becomes ever richer as the number system is extended.
    \item \textbf{Observer thread} (scale $\sigma \to$ sampling rate $R \to$ invisibility criterion). The categorical unification of the observer and periodic structure is completed in the `Periodic Categories' chapter.
\end{enumerate}

The three main threads ultimately converge on a core proposition: \textbf{in conservative systems, periodic structure may be closely related to conservation laws, and conservation laws are related to symmetry within the framework of Noether's theorem; but the true algebraic hinge of periodicity is the finite-order action, and conservation and symmetry are auxiliary tools for discovering, constraining and interpreting periods.} The task of periodic mathematics is, with group actions as its language and closed orbits as its criterion, to characterise precisely the deductive chain from finite-order action to periodicity, whilst clarifying the proper place of conservation and symmetry within it.

